% Standalone Overleaf source. Compile with pdfLaTeX; no external assets needed.
% This is an annotated skeleton, not a completed proposal or a novelty claim.
\documentclass[11pt,a4paper]{article}
\usepackage[T1]{fontenc}
\usepackage[utf8]{inputenc}
\usepackage[english]{babel}
\usepackage[margin=1in]{geometry}
\usepackage{amsmath,amssymb,amsthm}
\usepackage{xcolor}
\usepackage{enumitem}
\usepackage[expansion=false]{microtype}
\usepackage[hidelinks]{hyperref}
\setlist{itemsep=2pt,topsep=4pt}
\newcommand{\sg}{\operatorname{sg}}
\newcommand{\Tr}{\operatorname{Tr}}
\newcommand{\E}{\mathbb{E}}
\newcommand{\R}{\mathbb{R}}
\newcommand{\norm}[1]{\left\lVert #1\right\rVert}
\newif\ifdraftnotes
\draftnotestrue % Set to \draftnotesfalse to hide drafting prompts.
\newcommand{\draftnote}[1]{\ifdraftnotes\par\smallskip
  {\small\color{gray}\noindent\textit{To develop:} #1\par}\smallskip\fi}

\title{What Are We Optimizing When We Call \texttt{.detach()}?\\[0.4em]
\large Research proposal: stop-gradient dynamics, online adaptation, and forgetting}
\author{Amit Arad\\[0.3em]
\small Advisor: Assaf Shocher (Faculty of Data and Decision Sciences)}
\date{October 2026}

\begin{document}
\maketitle

\begin{abstract}
Stop-gradient and frozen reference models change how learning signals propagate.
This proposal studies the resulting dynamics as a model keeps adapting to a
changing input distribution. The emerging direction connects Lyapunov analysis
with online test-time adaptation and continual learning, using IT$^3$
(Idempotent Test-Time Training) as a concrete starting point. The immediate
goal is to specify a tractable student--reference system and ask when
adaptation remains stable, what it retains from earlier distributions, and
whether apparent forgetting can be reversed. PEIRA remains a mathematical
reference for convergence certificates; extending its exact-predictor model
is now an optional supporting direction. The models and theorem targets below
are proposals for discussion, not established results.
\end{abstract}

\section{Introduction: from stop-gradient to continual adaptation}
\label{sec:introduction}
\paragraph{Motivation.}
The numerical value of a loss does not determine the update implemented by a
computation graph. Detached outputs, frozen parameters, and moving references
can induce different training fields. In a stream of changing distributions,
we must distinguish stability of those dynamics from accuracy on the current
distribution and retention of earlier capabilities.

\paragraph{Provisional research question.}
Can an analysis of coupled adaptation dynamics inspired by IT$^3$ identify
conditions for stable adaptation and quantify the accompanying retention or
forgetting? Which conclusions depend on the backward rule, reference timescale,
and rate of distribution change?

\paragraph{Scope of the pivot.}
The aim is to develop a small, mathematically explicit problem for the research
group discussion. Online test-time adaptation supplies a concrete mechanism;
continual learning supplies questions about performance over a sequence of
distributions. These settings need not share the same objective. In particular,
the supplied IT$^3$ paper prioritizes current performance in its online setting
and permits deterioration on the original distribution~\cite{it3}.
Retention is an additional question for this proposal.
\draftnote{Choose whether the first result should address fixed-distribution
stability, tracking under drift, or retention after a switch. Do not promise
all three in the first theorem.}

\section{Stop-gradient as a training field}
\label{sec:framework}
\subsection{Forward values and backward rules}
In the forward pass, $\sg(u)=u$; automatic differentiation assigns the
stop-gradient node a zero Jacobian. Denote the resulting autodiff field by
$F(\theta)$, so a basic update and its continuous-time model are
\[
 \theta_{n+1}=\theta_n-\eta F(\theta_n),
 \qquad \dot\theta=-F(\theta).
\]
\draftnote{Include the scalar example where the forward loss and a potential
of $F$ differ, followed by the planar rotation example. Distinguish continuous
flow from discrete updates. Stop-gradient need not destroy conservativity.}

\subsection{Lift, differentiate, and restrict}
Replace each detached subexpression $\sg(g_k(\theta))$ by a shadow variable
$z_k$, giving $\widehat L(\theta,z)$ and $g=(g_1,\ldots,g_K)$. Then
\begin{equation}
 F(\theta)=\left.\nabla_\theta\widehat L(\theta,z)\right|_{z=g(\theta)},
 \qquad
 \nabla L_{\mathrm{forward}}=F+J_g^\top\nabla_z\widehat L.
 \label{eq:lift}
\end{equation}
For smooth functions,
\begin{equation}
 DF=\widehat L_{\theta\theta}+\widehat L_{\theta z}J_g.
 \label{eq:jacobian}
\end{equation}
On an open simply connected domain, a continuously differentiable field is
conservative precisely when $DF$ is symmetric (Poincaré).
\draftnote{Derive these identities and explain vector-valued shadows. Introduce
the game interpretation for suitable parameter blocks; a shadow is not
automatically an independent optimizing player.}

\subsection{Convergence, stability, and representation quality}
A candidate Lyapunov function $V$ satisfies
\[
 \frac{d}{dt}V(\theta(t))=-\langle\nabla V,F\rangle\leq 0.
\]
\draftnote{Distinguish a potential, a Lyapunov function, and an invariant.
State the additional boundedness, strictness, or analytic assumptions needed
for the particular convergence claim. An invariant alone does not exclude
cycles; convergence alone does not exclude collapse. Keep collapse stability
and Nash properties model-dependent.}

\section{Mathematical background: self-distillation and PEIRA}
\label{sec:peira}
\subsection{The motivating architectures}
\draftnote{Define views, encoder, predictor, target, and collapse. Explain the
shared stopped target in SimSiam and the EMA target in BYOL. Then introduce
the idealized model analyzed by PEIRA, distinguishing it from the full
architectures and their practical training procedures.}

\subsection{The exact-predictor setting}
For paired inputs $(x,y)$ and encoders $U(x),V(y)\in\R^k$, define
\begin{align}
 N&=\E[UU^\top+VV^\top], &
 \Sigma&=\E[VU^\top+UV^\top], &
 P^\star&=\Sigma(N+\lambda I)^{-1}.
\end{align}
PEIRA analyzes a population, function-space self-distillation flow with
this exact regularized predictor~\cite{peira}. Its certificate is
\begin{equation}
 \mathcal V_{\mathrm{SD}}(U,V)
 =\tfrac13\Tr[(N+\lambda I)^3]-\tfrac12\Tr(\Sigma^2).
\end{equation}
\draftnote{Write the directed regularized prediction losses and the exact
semi-gradient equations. State the spectral assumptions and explain nonlinear
CCA. Summarize convergence and equilibrium classification separately from
stability; the analyzed self-distillation system permits stable collapse.}

\subsection{PEIRA's new objective}
The method replaces prediction-error minimization by
\begin{equation}
 E_{\mathrm{PEIRA}}(U,V)
 =-\tfrac12\Tr[\Sigma(N+\lambda I)^{-1}]
   +\tfrac\lambda2\Tr N.
\end{equation}
\draftnote{Explain why this objective differs from the self-distillation
certificate. Use a canonical-mode calculation to show the threshold
$c>\lambda^2$, then state the full leading-subspace result with its assumptions.
Distinguish zero collapse from constant representations and clarify when
eligible nonconstant modes exist.}

\subsection{The boundary of the results}
\draftnote{Separate three issues: finite-speed predictor optimization; EMA
teacher dynamics; neural-network parameterization and stochastic training.
PEIRA leaves teacher-EMA convergence and stability to future work (footnote 2).
Its practical auxiliary-gradient algorithm also uses EMA feature statistics;
this is distinct from an EMA teacher. The exact auxiliary-gradient identity
is useful for studying when detached computations still implement an objective.}

\section{\texorpdfstring{IT$^3$}{IT3} as a concrete bridge}
\label{sec:it3}
\subsection{Prediction, reference, and the online extension}
In the supplied version of IT$^3$, the model accepts an input $x$ and an
additional prediction or neutral signal. Write
\[
 q_\theta(x)=f_\theta(x,0),\qquad T_\phi(x,q)=f_\phi(x,q).
\]
The test-time discrepancy compares $T_\phi(x,q_\theta(x))$ with $q_\theta(x)$.
For episodic adaptation, the reference weights are pretrained and fixed;
the adaptive model is reset between episodes. The online extension carries
the adaptive weights forward and uses an EMA reference~\cite[Sections 3.2--3.3]{it3}.
This motivates studying coupled dynamics and changing data.

\subsection{Frozen parameters are not detached outputs}
For analysis, consider the \emph{squared Euclidean} surrogate
\begin{equation}
 D(\theta,\phi;\rho)
 =\tfrac12\E_{x\sim\rho}
 \bigl\|T_\phi(x,q_\theta(x))-q_\theta(x)\bigr\|^2.
 \label{eq:it3surrogate}
\end{equation}
This is a proposed smooth surrogate; the supplied paper writes a norm
discrepancy. Put $r=T_\phi(x,q_\theta(x))-q_\theta(x)$ and
$J_q=\partial q_\theta(x)/\partial\theta$. Holding $\phi$ fixed during
differentiation gives
\begin{equation}
 G_{\rm frozen}(\theta,\phi;\rho)
 =\E_\rho\!\left[J_q^\top(\partial_q T_\phi-I)^\top r\right].
 \label{eq:frozen}
\end{equation}
Freezing reference parameters still allows the input derivative
$\partial_q T_\phi$ to transmit a learning signal.
Explicitly detaching the whole reference output instead gives
\begin{equation}
 G_{\rm detached}(\theta,\phi;\rho)=-\E_\rho[J_q^\top r].
 \label{eq:detached}
\end{equation}
The two expressions share forward values but generally differ as updates.
The detached-output version is a proposed comparison, not a claim about
IT$^3$'s implementation. Check the actual backward paths and adapted parameters
in the chosen code version before making an implementation claim.
Idempotence alone does not establish accuracy, orthogonal projection, or
preservation of old tasks.

\section{Proposed research: stability and retention under adaptation}
\label{sec:research}
\paragraph{Scope and status.}
The following model and research targets are working choices motivated by
the pivot. No convergence theorem, anti-forgetting guarantee, or novelty
claim is asserted.

\subsection{First task: specify a minimal coupled system}
Start with a finite-dimensional smooth model and a distribution $\rho_t$:
\begin{equation}
 \dot\theta=-G(\theta,\phi;\rho_t),\qquad
 \tau\dot\phi=\theta-\phi,\qquad \tau>0.
 \label{eq:teacher}
\end{equation}
Choose $G$ explicitly from \eqref{eq:frozen} or \eqref{eq:detached}.
Treat the ODE as an idealization, and separately specify the discrete
optimizer and EMA update. Begin with fixed $\rho$, then a two-distribution
switch, before considering smooth drift or stochastic training.
\draftnote{Choose a scalar or low-dimensional affine model whose shared
parameters couple predictions on old and new inputs. Avoid a model in
which each input can adapt independently. State the adapted parameters
and the pretrained initialization.}

\subsection{Second task: derive a certificate rather than assume one}
For the frozen-reference gradient field, with $\rho$ fixed and $\phi$ held
constant, the surrogate obeys $\dot D=-\|\nabla_\theta D\|^2$.
With an EMA reference, the same calculation instead yields
\begin{equation}
 \dot D=-\|\nabla_\theta D\|^2
 +\tau^{-1}\langle\nabla_\phi D,\theta-\phi\rangle.
 \label{eq:emaenergy}
\end{equation}
The added term has no automatic sign. If $\rho_t$ varies differentiably and
differentiation under the expectation is justified, a further term
$\partial_tD$ accounts for distribution change. Abrupt switches require
separate accounting for the jump in the evaluated objective.
An initial candidate to investigate is
\begin{equation}
 W(\theta,\phi;\rho)
 =D(\theta,\phi;\rho)+a\|\theta-\phi\|^2,\qquad a>0.
 \label{eq:composite}
\end{equation}
This is an ansatz, not an established Lyapunov function. For the detached
field, even the first term in the derivative must be recomputed as
$-\langle\nabla_\theta D,G_{\rm detached}\rangle$.
\draftnote{Compute equilibria and the full Jacobian in the chosen toy model.
Determine whether a local certificate, a tracking bound, or a counterexample
is the most informative first result. Check boundedness and non-isolated
equilibria before claiming convergence to a point.}

\subsection{Third task: measure retention separately from consistency}
For an earlier labeled evaluation distribution $\mu_A$, define
\begin{equation}
 R_A(\theta)=\E_{(x,y)\sim\mu_A}
 \ell_{\rm task}(q_\theta(x),y),\qquad
 \Delta_A(t)=R_A(\theta_t)-R_A(\theta_{t_A}),
 \label{eq:retention}
\end{equation}
where $t_A$ is a fixed reference time after learning or adapting to $A$.
Positive $\Delta_A$ means deterioration relative to that reference.
These labels are used for evaluation, not for test-time updates.
For smooth $R_A$,
\begin{equation}
 \frac{d}{dt}R_A(\theta_t)
 =-\langle\nabla_\theta R_A,G(\theta,\phi;\rho_t)\rangle.
 \label{eq:interference}
\end{equation}
Decreasing the adaptation discrepancy does not by itself establish retention:
an additional relation between its update and the old-task risk is required.

Use an $A\to B\to A$ sequence to distinguish immediate performance loss
from recovery after returning to an earlier distribution. Record recovery
speed and the amount of new adaptation required; recovery after retraining
alone does not prove that the original representation remained intact.
\draftnote{Use the forgetting readings to define recoverability. Decide
whether the first analysis concerns old-task risk, representation change,
or access to retained information, and justify the connection between
the chosen quantities.}

\subsection{Candidate outcomes for the first study}
\begin{enumerate}
 \item A local stability calculation or certificate for one explicit
 student--reference system, with conditions on $\tau$ and the backward rule.
 \item A tracking estimate under bounded distribution drift, stating both
 the regularity assumptions and what quantity is being tracked.
 \item A counterexample in which adaptation discrepancy falls while
 earlier-task error grows, identifying the missing retention assumption.
\end{enumerate}
These are alternative initial targets. A retention-preserving modification
should follow from the diagnosed mechanism and a stated information budget.
It must not silently use old labels or examples unavailable to the learner.

\subsection{Role of the earlier finite-speed predictor direction}
PEIRA remains useful for learning how to construct and interpret a
certificate. The earlier extension
\[
 \dot\theta=-F(\theta,p),\qquad
 \varepsilon\dot p=-\nabla_p R(\theta,p)
\]
can serve as a secondary theoretical exercise in tracking an auxiliary
optimum. It is no longer the primary application target. Any transfer of
its arguments to IT$^3$ must identify matching assumptions; a predictor
optimum and an EMA reference obey different dynamics.

\section{Reading priorities, experiments, and milestones}
\label{sec:plan}
\paragraph{Reading priorities.}
First, reconstruct IT$^3$'s training and adaptation graphs, especially the
frozen-reference and online EMA variants in the supplied arXiv
v1~\cite{it3}. Next, use the forgetting and continual-learning readings to
fix the learning setting and evaluation criterion. Return to PEIRA and
the stop-gradient convergence paper to identify applicable proof techniques.
\draftnote{Identify the exact forgetting and continual-learning position
papers in the reading discussion before adding bibliographic claims.
Read the supplied stop-gradient convergence and LyapuFlow papers for their
precise assumptions; a feedback certificate for a generative flow is not
automatically a training-dynamics certificate.}

\begin{enumerate}
 \item \textbf{Graph and model audit.} Specify forward expressions, trainable
 parameters, backward paths, and reference updates. Reproduce the derivatives
 of the selected squared surrogate.
 \item \textbf{Fixed-distribution baseline.} Locate equilibria, compute the
 coupled Jacobian, and compare frozen and detached fields under matched
 initialization.
 \item \textbf{Controlled switches.} Evaluate $A\to B\to A$ and gradual
 shifts while varying EMA speed, adaptation steps, and shift magnitude.
 Compare no adaptation, episodic reset, fixed reference with carried
 adaptive weights, and EMA reference with carried adaptive weights.
 \item \textbf{Diagnostics.} Track adaptation discrepancy, evaluation risk
 on current and earlier distributions, student--reference distance,
 representation change, and recovery time. Check candidate-certificate
 derivatives and repeat with smaller steps to detect numerical artifacts.
 \item \textbf{Meeting deliverable.} Bring one explicit model, its derived
 update equations, one informative phase portrait or switching experiment,
 and a sharply stated theorem question or counterexample.
\end{enumerate}
\draftnote{After the toy result is understood, select one small neural
experiment and audit related work to identify what is new. Record memory,
data-access, and compute budgets so comparisons answer the same question.}

\section{Expected contribution and limits}
The intended contribution is a precise connection between a backward rule,
a moving reference, and adaptation dynamics in a changing environment.
A useful initial result may explain a stability boundary, bound tracking
error, or show why a convergence certificate cannot also certify retention.
The proposal does not assume that consistency implies correctness, that EMA
prevents forgetting, or that every loss of old-task accuracy is irreversible.
\draftnote{Replace these intended outcomes with evidence after the first
analysis. Separate background results, reproduced calculations, and new
findings.}

\medskip
\noindent{\small\textit{*Acknowledgments.*} This proposal has developed through ongoing discussion and collaboration with Shoham Grunblat.}

\begin{thebibliography}{9}
\bibitem{it3}
N. Durasov, A. Shocher, D. Oner, G. Chechik, A. A. Efros, and P. Fua.
\emph{IT$^3$: Idempotent Test-Time Training}.
\href{https://arxiv.org/abs/2410.04201}{arXiv:2410.04201}.
Supplied version: v1, 5 October 2024; Sections 3.2--3.3.

\bibitem{peira}
M. Arbel, B. Terver, and J. Ponce.
\emph{PEIRA: Learning Predictive Encoders through Inter-View Regressor Alignment}.
arXiv:2605.17671, 2026. Supplied version: 17 May 2026.
\bibitem{games}
A. Letcher et al.
\emph{Differentiable Game Mechanics}.
Journal of Machine Learning Research, 20(84):1--40, 2019.
% TODO: Verify and add the primary BYOL, SimSiam, and Tian--Chen--Ganguli
% references, relevant singular-perturbation results, and the precursor
% self-distillation analysis cited by PEIRA. Do not infer novelty from this
% deliberately incomplete starter bibliography.
\end{thebibliography}

\appendix
\section{Derivations and audit of the original proposal}
\draftnote{Place the lift-Jacobian derivation, toy phase portraits, the frozen-reference and detached-output derivative calculations,
and toy-model stability derivations here. Keep the main proposal centered
on adaptation, retention, and the selected theorem question.}

\end{document}

